
\documentclass[12pt]{article}

% ---------- Packages ----------
\usepackage{amsmath}
\usepackage{amssymb}
\usepackage{amsthm}
\usepackage{geometry}
\usepackage{enumitem}
\usepackage{hyperref}

\geometry{a4paper, margin=1in}
\raggedbottom

% Allow consecutive empty headings to break across pages.
\AddToHook{cmd/section/before}{\par\penalty0}
\AddToHook{cmd/subsection/before}{\par\penalty0}

% ---------- Theorem Environments ----------
\newtheorem{definition}{Definition}
\newtheorem{theorem}{Theorem}
\newtheorem{lemma}{Lemma}
\newtheorem{corollary}{Corollary}
\newtheorem{example}{Example}

% ---------- Document ----------
\begin{document}

\title{\textbf{Probability and Statistics}}
\author{Annukul Ghale}


\maketitle

\tableofcontents
\newpage


% ==================================================
% CHAPTERS
% ==================================================

\section{Chapter 1 - Samples, Spaces and Events}
\subsection{Definitions}
\begin{definition}
The sample space is the set of all possible outcomes for the experiment. It is denoted by $S$ (or $\Omega$).
\end{definition}

\begin{definition}
An event is a subset of the sample space. The event occurs if
the actual outcome is an element of this subset.
\end{definition}

\begin{definition}
An event E is a simple event (or elementary event) if it
consists of a single element of the sample space S.
\end{definition}

\subsection{Notes}
The size, or cardinality, of a finite set is the number of elements it contains.

\begin{itemize}
\item The set difference $A \setminus B$ (``$A$ take away $B$'') is the set of elements in $A$ but not in $B$:
\[
A \setminus B = \{x : x \in A \text{ and } x \notin B\}.
\]
\item The symmetric difference $A \triangle B$ is the set of elements in either $A$ or $B$, but not both:
\[
A \triangle B = (A \setminus B) \cup (B \setminus A).
\]
\end{itemize}

\begin{itemize}
\item The complement of a set $A$ is the set of elements in the sample space that are not in $A$. It can be denoted by $A'$ or $A^c$:
\[
A^c = S \setminus A.
\]
\item Two sets $A$ and $B$ are disjoint (or mutually exclusive) if they have no elements in common:
\[
A \cap B = \varnothing.
\]
\end{itemize}

De Morgan's laws describe how complements interact with unions and intersections:
\[
(A \cup B)^c = A^c \cap B^c,
\qquad
(A \cap B)^c = A^c \cup B^c.
\]
In words, the complement of a union is the intersection of the complements, and the
complement of an intersection is the union of the complements.

\section{Chapter 2}
\subsection{Definitions}
\subsection{Notes}

\section{Chapter 3}
\subsection{Definitions}
\subsection{Notes}

\section{Chapter 4}
\subsection{Definitions}
\subsection{Notes}

\section{Chapter 5}
\subsection{Definitions}
\subsection{Notes}

\section{Chapter 6}
\subsection{Definitions}
\subsection{Notes}

\section{Chapter 7}
\subsection{Definitions}
\subsection{Notes}

\section{Chapter 8}
\subsection{Definitions}
\subsection{Notes}

\section{Chapter 9}
\subsection{Definitions}
\subsection{Notes}

\section{Chapter 10}
\subsection{Definitions}
\subsection{Notes}

\section{Chapter 11}
\subsection{Definitions}
\subsection{Notes}

\section{Chapter 12}
\subsection{Definitions}
\subsection{Notes}

\section{Chapter 13}
\subsection{Definitions}
\subsection{Notes}

\section{Chapter 14}
\subsection{Definitions}
\subsection{Notes}

\section{Chapter 15}
\subsection{Definitions}
\subsection{Notes}

\section{Chapter 16}
\subsection{Definitions}
\subsection{Notes}

\section{Chapter 17}
\subsection{Definitions}
\subsection{Notes}

\section{Chapter 18}
\subsection{Definitions}
\subsection{Notes}

\section{Chapter 19}
\subsection{Definitions}
\subsection{Notes}

\section{Chapter 20}
\subsection{Definitions}
\subsection{Notes}




\end{document}
